% Zdroj: PDF priklady-jaro-2024.pdf, fyzická str. 1-3. Značka: společné okruhy. Zařazeno: I4.
\section{Ovladač pro řadič disku}
\szzotazka{jaro 2024}{2}{Ovladač pro řadič disku}

\noindent\textbf{Zadání.}\par
Napište implementaci funkce ovladače disku pro načtení jednoho bloku.

Uvažovaný disk je řízen pěti paměťově mapovanými registry, které slouží k~zápisu příkazů
(a~jejich parametrů) a~ke čtení aktuálního stavu zařízení (uvažujeme zařízení bez podpory
přerušení, s~obsluhou s~aktivním čekáním).

\begin{center}
\begin{tabular}{rllp{0.6\linewidth}}
\toprule
Offset & Registr & Typ & Popis registru \\
\midrule
0  & Status  & R & Bitové pole pro aktuální stav řadiče (pouze pro čtení, zápisy jsou ignorovány). \\
4  & Size    & R & Velikost disku v~blocích (pouze pro čtení, zápisy jsou ignorovány). \\
8  & Command & W & Zápis do registru spustí zadaný příkaz (1 čtení, 2 zápis). \\
12 & LBA     & W & Logická adresa bloku pro příkaz zadaný v~\textit{Command}. \\
16 & DMA     & W & Fyzická adresa paměti pro uložení načtených dat z~disku. \\
\bottomrule
\end{tabular}
\end{center}

\noindent Registr Status (offset 0) je bitové pole:
\begin{center}
\begin{tabular}{|c|c|c|c|c|c|}
\hline
31 & 29 & $\dots$ & 2 & 1 & 0 \\
\hline
0 & 0 & $\dots$ & 0 & BUSY & ERR \\
\hline
\end{tabular}
\end{center}
\noindent BUSY -- Bit je nastaven po zadání příkazu a~vynulován po dokončení operace.\par
\noindent ERR -- Bit signalizující chybový stav (1 znamená chybu).

\medskip
Bloky jsou na disku adresovány obvyklým způsobem pomocí LBA, jednotlivé bloky mají pevně
danou velikost 512 bajtů. Přenos čtených či zapisovaných dat probíhá pomocí DMA (tedy
z~pohledu této otázky není třeba řešit nic víc než nastavení fyzické adresy pro uložení
dat, vlastní čtení z~disku a~zápis do paměti je plně v~režii řadiče), dokončení operace je
signalizováno stavovým registrem (předpokládáme 32-bitový systém, velikost disku je
omezena na 2\,TB).

Dopište těla dvou níže uvedených funkcí a~navrhněte strukturu pro uchování informací
o~disku (předpokládáme, že k~systému může být připojeno více disků a~ty jsou na úrovni
jádra operačního systému rozlišeny různými instancemi struktury \texttt{disk\_t}). Obě
funkce vrací \texttt{true} pokud se daná operace zdařila, jinak \texttt{false} (bez
dalších detailů).

Funkce \texttt{disk\_init} je volána jednou při inicializaci daného zařízení, parametr
\texttt{register\_address} je virtuální adresa mapovaných registrů (tj. přímo adresa
stavového registru). Funkce \texttt{disk\_read\_block\_waiting} přečte jeden blok z~disku.
Načtený blok má být zapsán na fyzickou adresu \texttt{data\_phys\_addr} (váš kód může
předpokládat její správnost). Návrat z~funkce musí proběhnout až po dokončení čtení (je
možné využít aktivní čekání).

Vaše implementace musí alespoň triviálním způsobem ošetřit možný současný přístup k~disku
z~více procesů.

\begin{lstlisting}[language=C]
typedef struct { ... } disk_t;

bool disk_init(disk_t *disk, uint32_t register_address) { ... }

bool disk_read_block_waiting(disk_t *disk, size_t lba, uint32_t data_phys_addr) { ... }
\end{lstlisting}

\medskip
\noindent\textbf{Náčrt řešení.}\par
Nástin zdrojového kódu (bez symbolických konstant a~komentářů):

\begin{lstlisting}[language=C]
typedef volatile struct {
    uint32_t status;
    uint32_t size;
    uint32_t command;
    uint32_t lba;
    uint32_t dma;
} disk_regs_t;

typedef struct {
    size_t max_lba;
    disk_regs_t *ctl;
    mutex_t mutex;
} disk_t;

static bool disk_is_ok(disk_t *disk) { return disk->ctl->status & 1 == 0; }
static bool disk_is_ready(disk_t *disk) { return disk->ctl->status & 2 == 0; }
static bool disk_wait_for_ready(disk_t *disk) {
    while (!disk_is_ready(disk)) {
        if (!disk_is_ok(disk)) {
            return false;
        }
    }
    return true;
}

bool disk_init(disk_t *disk, uint32_t register_address) {
    disk->ctl = (disk_regs_t *) register_address;
    disk->max_lba = disk->ctl->size;
    mutex_init(&disk->mutex);
    return disk_wait_for_ready(disk);
}

bool disk_read_block_waiting(disk_t *disk, size_t lba, uint32_t data_phys_addr) {
    if (lba >= disk->max_lba) {
        return false;
    }
    mutex_lock(&disk->mutex);
    bool ok = disk_wait_for_ready(disk);
    if (!ok) {
        mutex_unlock(&disk->mutex);
        return false;
    }

    disk->ctl->lba = lba;
    disk->ctl->dma = data_phys_addr;
    disk->ctl->command = 1;

    ok = disk_wait_for_ready(disk);
    mutex_unlock(&disk->mutex);

    return ok;
}
\end{lstlisting}

\medskip
\noindent\textit{Doplňující vysvětlení (nad rámec oficiálního náčrtu):} Ve funkcích
\texttt{disk\_is\_ok} a~\texttt{disk\_is\_ready} chybí v~náčrtu závorky. V~C má \texttt{==} vyšší
prioritu než bitové \texttt{\&}, takže \texttt{status \& 1 == 0} znamená \texttt{status \& (1 == 0)},
tj.\ \texttt{status \& 0}, a~to je vždy~0. Správně má být \texttt{(status \& 1) == 0}
a~\texttt{(status \& 2) == 0}. Zbytek náčrtu: struktura registrů je \texttt{volatile}, aby překladač
čtení stavu v~čekací smyčce pokaždé skutečně provedl a~zápisy parametrů a~příkazu nepřeházel.
Mutex serializuje přístup více procesů k~jediné sadě registrů a~odemyká se na všech návratových
cestách. Parametry (LBA, DMA) se zapisují \emph{před} příkazem, protože zápis do \texttt{Command}
operaci rovnou spouští.
